\chapter{Projekt od zera: budujemy maly serwis krok po kroku}

\section{Cel projektu}
Na koniec przejdziemy przez caly, realistyczny mini-projekt. Zbudujemy prosty serwis
zlozony z:

\begin{itemize}
  \item strony glownej,
  \item listy wpisow,
  \item wspolnych metadata,
  \item funkcji pomocniczych,
  \item danych w pliku JSON,
  \item kilku importowanych chunkow.
\end{itemize}

Naszym celem nie jest stworzenie ogromnej aplikacji, tylko pokazanie, jak elementy
MoonChunk lacza sie w sensowna calosc.

\section{Struktura plikow projektu}
Przyjmijmy taki uklad:

\begin{lstlisting}[style=basecode,caption={Struktura przykladowego projektu}]
book-demo/
  content/
    site.mncnk
    metadata.mncnk
    blog.mncnk
    sections.mncnk
    data/
      posts.json
  run.js
\end{lstlisting}

\section{Dane wpisow}
Zacznijmy od danych. Plik \mc{posts.json} moze wygladac tak:

\begin{lstlisting}[style=basecode,caption={Przykladowy plik posts.json}]
[
  {
    "title": "MoonChunk od podstaw",
    "slug": "moonchunk-od-podstaw",
    "description": "Pierwszy kontakt z jezykiem i podstawowa skladnia.",
    "published": true,
    "views": 120
  },
  {
    "title": "Funkcje i petle",
    "slug": "funkcje-i-petle",
    "description": "Jak przetwarzac wiele elementow i porzadkowac logike.",
    "published": true,
    "views": 90
  },
  {
    "title": "Wpis roboczy",
    "slug": "wpis-roboczy",
    "description": "Jeszcze nie powinien byc pokazany publicznie.",
    "published": false,
    "views": 5
  }
]
\end{lstlisting}

\section{Wspolne metadata}
Nastepnie wydzielmy wspolna konfiguracje strony do osobnego chunku.

\begin{lstlisting}[style=moonchunk,caption={Plik metadata.mncnk}]
export chunk "Metadata" {
  lang: "pl";
  charset: "utf-8";
  viewport: "width=device-width, initial-scale=1";
  metaAuthor: "MoonChunk Team";
  metaRobots: "index,follow";
  themeColor: "#111827";
  styles: "<link rel='stylesheet' href='/assets/app.css'>";
  footer: "<footer><small>MoonChunk Demo</small></footer>";
};
\end{lstlisting}

To bardzo praktyczny krok. Dzieki temu nie musimy powtarzac tych samych pol w wielu
chunkach.

\section{Wspolna sekcja powitalna}
Zrobmy tez dodatkowy chunk z sekcja powitalna.

\begin{lstlisting}[style=moonchunk,caption={Plik sections.mncnk}]
export chunk "HeroSection" {
  page "" {
    content {
      <section class="hero">
        <h1>MoonChunk Blog</h1>
        <p>Statyczny serwis budowany z danych i logiki.</p>
      </section>
    };
  };
};
\end{lstlisting}

W tym miejscu warto zwrocic uwage na ciekawy aspekt: chunk moze zawierac w sobie
\mc{page}. Oznacza to, ze logika strony moze byc skladana z mniejszych czesci.

\section{Logika bloga}
Teraz utworzmy osobny plik odpowiedzialny za generowanie strony z wpisami.

\begin{lstlisting}[style=moonchunk,caption={Plik blog.mncnk}]
export chunk "BlogPages" {
  const posts = data("./data/posts.json");

  function badge(views: int): string {
    return views > 100 ? "Popularny" : "Nowy";
  }

  page "blog" {
    content {
      <section>
        <h2>Lista wpisow</h2>
      </section>
    };

    for (let i: int = 0; i < 3; i++) {
      if (posts[i].published) {
        content {
          <article class="post-card">
            <h3>{posts[i].title}</h3>
            <p>{posts[i].description}</p>
            <p>Status: {badge(posts[i].views)}</p>
            <p>Views: {posts[i].views}</p>
            <a href={"/posts/" + posts[i].slug}>Czytaj wiecej</a>
          </article>
        };
      };
    };
  };
};
\end{lstlisting}

W tym fragmencie mamy od razu kilka istotnych elementow:

\begin{itemize}
  \item wczytanie danych przez \mc{data(...)},
  \item funkcje pomocnicza,
  \item petle \mc{for},
  \item warunek \mc{if},
  \item dynamiczne atrybuty HTML,
  \item generowanie wielu podobnych kart.
\end{itemize}

\section{Plik glowny}
Na koniec skladamy wszystko w punkcie wejscia.

\begin{lstlisting}[style=moonchunk,caption={Plik site.mncnk}]
import { Metadata } from "./metadata.mncnk";
import { BlogPages } from "./blog.mncnk";
import * as Sections from "./sections.mncnk";

chunk "Main" {
  @include Metadata;
  @include BlogPages;
  output: "./dist";
  title: "MoonChunk Blog";
  metaDescription: "Przykladowy serwis statyczny zbudowany w MoonChunk.";

  page "" {
    content {
      <section>
        <h1>Start</h1>
        <p>To jest strona glowna serwisu.</p>
      </section>
    };
  };
};

moon(Main);
\end{lstlisting}

\section{Co zrobi ten projekt}
Po uruchomieniu projekt wygeneruje co najmniej:

\begin{itemize}
  \item strone glowna \mc{index.html},
  \item strone bloga \mc{blog.html},
  \item wspolny dokument z odpowiednimi metadata,
  \item HTML z kartami dla wpisow, ktore maja \mc{published = true}.
\end{itemize}

Dzieki temu powstaje maly, ale realistyczny serwis statyczny.

\section{Minimalny runner do tego projektu}
\begin{lstlisting}[style=basecode,caption={Plik run.js dla projektu}]
const { executeMoonChunkFile } = require("moonchunk");

const result = executeMoonChunkFile("./content/site.mncnk");

if (!result.ok) {
  console.error(result.diagnostics);
  process.exit(1);
}

console.log(result.generatedFiles || []);
\end{lstlisting}

To wszystko. Od strony uzytkownika tyle wystarczy, aby uruchomic generator.

\section{Jak rozwijac ten projekt dalej}
To dopiero poczatek. Z tego miejsca mozna isc w wielu kierunkach:

\begin{itemize}
  \item dodac osobne strony \mc{posts/...},
  \item rozbudowac metadata dla Open Graph i Twittera,
  \item dodac wiecej funkcji pomocniczych,
  \item wprowadzic wspolne chunki dla sekcji strony,
  \item podzielic dane na kilka plikow JSON,
  \item wzbogacic layout przez CSS i skrypty dolaczane w metadata.
\end{itemize}

\section{Czego ten projekt uczy}
Ten projekt jest wartosciowy nie dlatego, ze jest wielki, tylko dlatego, ze przechodzi
przez niemal wszystkie fundamenty MoonChunk:

\begin{itemize}
  \item strukture plikow,
  \item importy,
  \item eksportowane chunki,
  \item \mc{@include},
  \item \mc{moon(...)},
  \item metadata,
  \item funkcje,
  \item petle i warunki,
  \item dane JSON,
  \item generowanie HTML.
\end{itemize}

Jesli rozumiesz ten przyklad od poczatku do konca, masz bardzo solidne podstawy do
budowania wlasnych projektow w MoonChunk.

\section{Ostatnia rada praktyczna}
Kiedy tworzysz pierwszy prawdziwy projekt, nie probuj od razu zbudowac wszystkiego.
Najpierw uruchom jedna strone, potem dodaj wspolne metadata, potem dane, potem petle,
a dopiero na koncu bardziej zaawansowane elementy. To najspokojniejsza droga do
opanowania jezyka.
